\section{Unanticipated Recap}

\begin{quotex}
Nature does not make men equal; souls differ still more than bodies. \flright{\textsc{Antoine Fabre d'Olivet}}

\end{quotex}
As I've been pondering which of two alternative endings to use for Post \#1001 – if it ever appears at all – some recent inquiries have prompted this recapitulation before moving on to other topics. Perhaps a generation from now, some PhD student will write a dissertation of futile reactionary movements of the early XXIst century; he may then follow each obscure reference and come to a conclusion. Esoteric writings always leave some things unsaid.

It is pointless to explain a joke, since it then loses its purpose as a joke: it can no longer be funny. But is it pointless to explain a series of posts? After all, if someone cannot connect the dots, perhaps it could be my own fault. One of my creative writing teachers was highly critical of me. Since he preferred the vapid doodlings of the women in the class, I never took him seriously. Maybe he was correct. Is an inscrutable metaphor still a metaphor?

\paragraph{Living a Worldview}
Some seem to believe that we are an advocacy site, comparable to movements that retain a superficial resemblance. But advocacy is the role of a sophist, not a philosopher who strives to see things objectively and dispassionately. The latter has justice as the goal while the former strives for victory of some sort.

Unfortunately, unless the situation is understood objectively, the results of actions will be mixed and victory ephemeral. At that point, another virtue is necessary: courage or fortitude. A man should live by his convictions. That is the essence of the point made by C S Peirce.

\paragraph{Life Aim}
The World offers a myriad of options for living one's life. Some will claim they ``will decide for themselves", but based on what? Where do their thoughts come from? How can they understand possibilities of manifestation that have never occurred to them? Walk through a large university library; there are hundreds of thousands of volumes with conflicting facts and opinions. Can anyone master that?

Hence, a worldview can be no more than arbitrary no matter how fervently held. Nevertheless, there may be a way out: a thought from a higher source may intrude. If your mind is attuned to it, you may be able to follow it to escape the chaos and randomness of world events. Such men who are committed to that course have a life aim. Otherwise, for everyone else, life is little more than Brownian motion.

\paragraph{Blinded by Science}
Some still hold that science will open up that path. Yet science is always in the process of becoming, it is incomplete and subject to change, so how can the scientific method show us the path? So this is the first test: who can live solely by science under those circumstances? I will now provide some bullet points summarizing ideas from recent posts. You can follow along. The path will gradually be made clear, or at least, you can see where we have gone wrong. On the other hand, you may see where certain men have lacked the courage to follow up on their world view.

\begin{itemize}
\item \textbf{Natural Selection}. All life arose from random variation and natural selection. But this does not explain why there are hippopotamuses and not unicorns. It does not explain the observed sequence to more complex life forms. At best, it can only be an alternative theory to design, not a refutation. Otherwise, there would be no talk of living in a simulation, etc. Natural selection is accepted to avoid facing up to the possibility of a designer. 
\item \textbf{Animality}. The corollary is that human beings are just a species of animal. Does anyone actually advocate living like animals, which are motivated strictly by fear, sex, and hunger, or attraction and repulsion? Few will deny the efficacy of consciousness, rationality, or morality, but on what grounds? Animals are not bothered by moral considerations. 
\item \textbf{Rationality}. The average IQ of physicists is much higher than that of biologist, so we consider the views of Roger Penrose. He correctly understands that, based on Gödel's theorem, rationality cannot be explained as a biochemical neurological process. Unfortunately, he then assumes that some future theory of quantum gravity will somehow explain consciousness and rationality. However, that is a technique of evasion. 
\item \textbf{Morality}. Scientists don't think very deeply about things, since they are only interested in quantitative results which are beyond – or actually, below – thought. However, the philosopher Thomas Nagel sees clearly that natural selection cannot explain the qualities that make us human. Hence, he postulates an environment that will give rise to such qualities, although it cannot be known by science. Again, Nagel refuses to live by the results of his insights. 
\item \textbf{Scientific closure}. Hence, what matters most cannot be explained by science, that is, there is no epistemological scientific closure: something else must be added. The scholastic principle is that the real is intelligible, yet science does not account for the intelligibility of the world. There is no scientific beginning, since physics can never get to time=0 of the Big Bang. 
\item \textbf{Degrees of Existence}. So if the material explanation of the world is incomplete, there must be higher, non-material degrees of existence. Wolfgang Smith shows that the qualitative corporeal world is not the quantitative physical world. Moreover, there is vertical causation, so that the higher degrees of existence form the lower levels. A fortiori, it is not matter that creates the higher degrees. Smith shows that a corporeal body can be freed from its physical counterpart. Although it may look identical, it will have different powers. Specifically, a ``resurrection body" can be explained in that way. 
\end{itemize}
This is the rational path, the path of the intellect. It is more a question of will rather than knowledge. If the world is not epistemologically closes, then there is the possibility of revelation, or faith, intruding into the world from a higher state. If the physical world is just one state of being, then there is the possibility of super-physical events, or miracles. This is all rational, unless you want to stop at one of the bullet points. But which one?

\paragraph{Unobvious Conclusions}
So how can anyone take this as a ``criticism of secularism and in favor of a religiously and ritually structured society"? Rather, it is ``in favor" of the eternally valid. What is the value in criticizing an illusion? There is only value in dispelling it. It is not a question of ``choosing" one view over another, as though they are on the same level.

If someone self-identifies as an animal, what value is there in debating him? Does a man argue with his dog?

So here is a young man rather upset with an atheist's support of Hillary Clinton: \textit{Sam Harris VS Milo Yiannopoulos on Trump}\footnote{\url{https://www.youtube.com/watch?v=0nOAT60NYDs}}. Sam Harris is a neuroscientist who claims that ``people's thoughts and intentions emerge from background causes of which we are unaware and over which we exert no conscious control." In other words, Harris' support of Hillary is not a rational choice but more likely the result of abnormal brain chemistry.

The attentive reader will note that Gornahoor agrees with that claim. However, Gornahoor offers what Harris cannot (no, psychedelics like MDMA are of no value): a way out, a practice to ``exert conscious control" over one's thoughts. If he were an astute and erudite student of religion, rather than a semi-educated polemicist, he may have noticed that as fundamental to those of spiritual depth.

Another misconception is that somehow we are appreciative of the so-called alt right. This would be a good test: hold a ``heavy metal" night at the Pulse nightclub in Orlando and see who shows up. There will be a large overlap with the alt right if anyone wants to do the Venn diagram. But we won't be there.



\flrightit{Posted on 2016-07-06 by Cologero }

\begin{center}* * *\end{center}

\begin{footnotesize}\begin{sffamily}



\texttt{Logres on 2016-07-18 at 00:24 said: }

``What is the value in criticizing an illusion? There is only value in dispelling it." I spent most of my college years in marshaling plans to work through BF Skinner, Karl Marx, etc, in order to reach rational conclusions that were educated. At some point, (and probably due greatly from happy influence), I grasped what my old college professor John Reist was saying about a book list of ``read before you die". It's also tiring to drudge in a world dominated by atheist worldviews (practically) and then rush home to stay up late reading the latest attack on consciousness, soul, etc. Their debates might be more interesting if we weren't living with the practical fruit of it all around us. Eg, the Huxley-Darwin debates were interesting, superficially, because they meant something, taking place in a world still largely dominated by Faith. I am not even really sure the new atheists believe, deeply believe, that there is no God, just that He is largely irrelevant. If there is no God, what possible good would holding a debate do anyone?


\end{sffamily}\end{footnotesize}
